% \documentclass{report}

% \usepackage{graphicx}
% \usepackage{dcolumn}
% \usepackage{bm}
% \usepackage{amsfonts, amssymb, amsmath}
% \usepackage{braket}

%\newcommand{\ket}[1]{\textstyle \left| #1 \right\rangle \displaystyle}
%\newcommand{\bra}[1]{\textstyle \left\langle #1 \right| \displaystyle}
%\newcommand{\braket}[2]{\textstyle \left\langle #1 \left|  #2 \right\rangle\right. \displaystyle}
%\newcommand{\im}{\mathrm{i}}
%\newcommand{\e}{\mathrm{e}}
%\newcommand{\pwrt}[2]{\frac{\partial #1}{\partial #2}}
%\newcommand{\diff}{\mathrm{d}}

%\begin{document}

\subsection{ Full symmetrized wavefunction, including envelope function
    anisotropy }
\label{section:full_wavefunctions}

This section is a collection of information derived in previous sections, and is
collected here to make referencing of the form of our wavefunctions for various
symmetries swift and convenient for users and maintainers of the software.

We have a few preliminaries to define. We may start with what we in this
document call the \emph{axial Bloch function}. In our system, the Bloch
functions we are most concerned are the six energy minima of the conduction band
of bulk silicon. Each lies along a particular crystallographic $k$-space axis,
which may be labeled $x$, $y$, and $z$. We denote the location in $k$-space where each minima is located as follows
\begin{equation}
    \bm{k}_{\pm q} = \pm \eta_{\text{Si}}k_{0}\hat{q},
\end{equation}
where
\begin{equation}
    k_{0} = \frac{2\pi}{a_{\text{Si}}}.
\end{equation}
Above $a_{\text{Si}} = \SI{.5431020511}{\nano\metre}$ is the lattice parameter of the silicon crystal, and $\eta_{\text{Si}} \approx .85$.

We use the following notation to describe one of these valley Bloch functions
\begin{align}
\Theta^{\left(\pm q\right)}
    &\coloneqq
        \e^{\im \bm{k}_{\pm q}\cdot\bm{r}}
        u_{\bm{k}_{\pm q}\bm{r}}
,\\
    &=
        \e^{\im \eta_{\text{Si}}k_{0}r\chi_{q}}
        u_{\bm{k}_{\pm q}\bm{r}}
.
\end{align}
In the previous section we found that our choice of envelope function
basis, along with certain parameters specifying functions within that basis,
allowed us to factor the envelope function from the Bloch functions of opposite
valleys. These factored out Bloch function summations have the following form
\begin{equation}
    \Theta^{\left(q\right)}_{\xi}
    =
        \Theta^{\left(+q\right)}
            +
        \xi
        \Theta^{\left(-q\right)}.
\end{equation}
We describe this form as the \emph{axial Bloch function}.

There is a feature of notation below we should discuss; previously, we defined
the $m$ index of functions in our basis to extend over the nonnegative integers.
Because some of the basis functions must be set to zero depending on the parity
of $m$, we introduce new notation to reflect this. Namely, $m_{e}$ extends over
the nonegative even integers, including zero. Whereas, $m_{o}$ should be
understood to extend over the nonnegative odd integers. When invoked without
subscript, $m$ retains its original domain of all nonnegative integers.

Now, as we build up our multi-valley wavefunctions we consider what we will
describe here as the \emph{$z$-longitudinal single-axis basis wavefunctions}.
These wavefunction terms are longitudinal in the sense that they only include
parts of the full wavefunction with only $z$-axial Bloch functions used in their
construction. We found that, for $\left(\kappa,i\right) = \left(A_{n},0\right)$,
$\left(E,i\right)$, and $\left(T_{n},z\right)$ the $z$-longitudinal single-axis
basis wavefunctions to go as
\begin{align}
    \psi_{\parallel_{0},m_{e}lk\sigma\omega}^{\left(z\right)}
        & \coloneqq
            \phi_{klm_{e}\zeta}^{(z)}
            \Theta^{\left(z\right)}
            _{\sigma\omega\left(-1\right)^{l+m_{e}/2}}
\tag{\ref{envelope:longitudinal:z:0}+\ref{single-axis:longitudinal:z:0}}
,\\
    \psi_{\parallel_{1},mlk\sigma\omega}^{\left(z\right)}
        & \coloneqq
            \left\{
                \phi_{klm\zeta_{0}}^{(x)}
                +
                \sigma\left(-1\right)^{l+m}\im^{+m}
                \phi_{klm\zeta_{1}}^{(y)}
            \right\}
            \Theta^{\left(z\right)}
            _{\omega\left(-1\right)^{m}}
\tag{\ref{envelope:longitudinal:z:1}+\ref{single-axis:longitudinal:z:1}}
.
\end{align}
These $z$-longitudinal functions were constructed from the \emph{naive envelope
basis function} denoted by lowercase $\phi$ with suitable azimuthal constants
$\zeta$ as given below. 
\begin{align}
    \zeta
        & =
            \sigma\left(-1\right)^{m_{e}/2}
,\\
    \zeta_{0}
        &=
            \left(-1\right)^{l}
\quad
    \zeta_{1}
        =
            \left(-1\right)^{l+m}
.
\end{align}

We have enough information now to construct the full wavefunctions for the
$\kappa = A,E$ symmetries. For those states, we will need $q$-longitudinal
single-axis wavefunctions for $q = x, y$ as well as $q = z$. Fortunately, we can
invoke \eqref{idx:env:xz} and \eqref{idx:env:yz} to find the single-axis
wavefunctions for the $x$ and $y$ axes. The relevant cyclic transformations go
simply as $\psi^{\left(x\right)} = \left(yzx\right)\psi^{\left(z\right)}$ and
$\psi^{\left(y\right)} = \left(zxy\right)\psi^{\left(z\right)}$. These
transformations straightforwardly swap the $z$ index in the above expressions
for the $x$ and $y$ expressions respectively. Following this line of reasoning
and using the equations mentioned above we can state the forms of the full $A$,
$E$ wavefunctions to be
\begin{align}
    \Psi_{\varrho,mlk}^{\left(A_{n}\right)}
        &=
            \frac{1}{3}
            \psi_{\varrho,mlk\sigma^{\left(A_{n}\right)}\left(+1\right)}
            ^{\left(z\right)}
                +
            \frac{1}{3}
            \psi_{\varrho,mlk\sigma^{\left(A_{n}\right)}\left(+1\right)}
            ^{\left(x\right)}
                +
            \frac{1}{3}
            \psi_{\varrho,mlk\sigma^{\left(A_{n}\right)}\left(+1\right)}
            ^{\left(y\right)}
\label{wavefunction:A}
,\\
    \Psi_{\varrho,mlk\sigma}^{\left(E,0\right)}
        &=
            \frac{1}{3}
            \psi_{\varrho,mlk\sigma\left(+1\right)}^{\left(z\right)}
                +
            \frac{\tau}{3}
            \psi_{\varrho,mlk\sigma\left(+1\right)}^{\left(x\right)}
                +
            \frac{\tau^{2}}{3}
            \psi_{\varrho,mlk\sigma\left(+1\right)}^{\left(y\right)}
\label{wavefunction:E:0}
,\\
    \Psi_{\varrho,mlk\sigma}^{\left(E,1\right)}
        &=
            \frac{1}{3}
            \psi_{\varrho,mlk\sigma\left(+1\right)}^{\left(z\right)}
                +
            \frac{\tau^{2}}{3}
            \psi_{\varrho,mlk\sigma\left(+1\right)}^{\left(x\right)}
                +
            \frac{\tau}{3}
            \psi_{\varrho,mlk\sigma\left(+1\right)}^{\left(y\right)}
\label{wavefunction:E:1}
.
\end{align}
Above $\tau = -\frac{1}{2} + \im\frac{\sqrt{3}}{2}$, one of the third roots of
unity. The $A$-type states are nondegenerate, so we drop the degeneracy index
above. We retain the degeneracy index for the doubly degenerate $E$-type states.
While $\sigma$ is a good quantum number for the $A$ and $T$ states, it is not
for the $E$ states. As such, since our basis respects this symmetry, we should
include contributions from both $\sigma=-1$ and $\sigma=+1$ when considering
$\kappa = E$. When necessary we can treat $\sigma$ as a basis index like $m$,
$l$, or $k$. Finally, we introduce a $\varrho = \parallel_{0}, \parallel_{1}$
index that we will describe here as the single-axis functional form index.

For the $A$ and $E$ symmetries, only $q$-longitudinal single valley terms are
necessary to construct the full multi-valley wavefunction. In essence the idea
of a "$z$-longitudinal" or "$z$-transverse" single-axis wavefunction is
irrelevant for these symmetries -- the three axes are all present in the
wavefunction, and are in some sense interchangeable. As we will see, however,
for $T$ symmetries this distinction between longitudinal and transverse in
relation to a given valley is relevant. In the previous section we were able to
find an expression for transverse forms for $\kappa = T_{n}$ states. The
$T$-states are triply degenerate, and each degeneracy may be valley indexed. If
we consider the $\left(\kappa, i\right) = \left(T_{n},q\right)$ state, we can
define $q$-longitudinal wavefunction terms to be those that include the
$q$-valley axial Bloch functions, while $q$-transverse wavefunction terms are
defined as those which include axial Bloch functions with valleys perpendicular
to $q$.

To wit, we found these \emph{$z$-transverse single-axis basis wavefunctions} to
be
\begin{align}
\Psi_{\perp_{0},m_{o}lk\zeta}^{\left(T_{n},z\right)}
    &=
        \phi_{m_{o}lk\zeta}^{\left(x\right)}
        \Theta^{\left(x\right)}_{\zeta\left(-1\right)^{l}}
        -
        \zeta\sigma^{\left(T_{n}\right)}\im^{+m}
        \phi_{m_{o}lk\bar{\zeta}}^{\left(y\right)}
        \Theta^{\left(y\right)}_{\zeta\left(-1\right)^{l}}        
,\\
\Psi_{\perp_{1},mlk}^{\left(T_{n},z\right)}
    &=
        % \phi_{mlk\left(-1\right)^{l+1}}^{\left(y\right)}
        \phi_{mlk\bar{\zeta}_{0}}^{\left(y\right)}
        \Theta^{\left(x\right)}_{\left(-1\right)^{m+1}}
        -
        \sigma^{\left(T_{n}\right)}\left(-1\right)^{l+m}\im^{+m}
        % \phi_{mlk\left(-1\right)^{l+m+1}}^{\left(x\right)}
        \phi_{mlk\bar{\zeta}_{1}}^{\left(x\right)}
        \Theta^{\left(y\right)}_{\left(-1\right)^{m+1}}        
,\\
\Psi_{\perp_{2},mlk}^{\left(T_{n},z\right)}
    &=
        \phi_{mlk\bar{\zeta}_{1}}^{\left(z\right)}
        \Theta^{\left(x\right)}_{\left(-1\right)^{m}}
        -
        \sigma^{\left(T_{n}\right)}\left(-1\right)^{l}\im^{+m}
        \phi_{mlk\bar{\zeta}_{0}}^{\left(z\right)}
        \Theta^{\left(y\right)}_{\left(-1\right)^{m}}
.
\end{align}
It should be made clear that in transverse forms, by definition, the valley
index of the Bloch functions functions are perpendicular to the valley index of
the degeneracy index. As we examine the naive envelope basis functions, it is
clear that various valley indices are used above and that which valley index is
used for the naive functions depends upon which form index $\perp_{0}$,
$\perp_{1}$, and $\perp_{2}$ is used. This is similar to the 


We can articulate longitudinal forms for the multi-valley wavefunction with
$\kappa = T_{n}$ built up from the \emph{$z$-longitudinal single-axis basis
wavefunction} we introduced earlier. Namely, for $\rho = \parallel_{0},
\parallel_{1}$ we find
\begin{equation}
    \Psi_{\varrho,mlk}^{\left(T_{n},z\right)}
        =
            \psi_
            {\varrho,mlk\sigma^{\left(T_{n}\right)}\left(-1\right)}
            ^{\left(z\right)}
.
\end{equation}

Above, we articulated wavefunctions for the states of $\left(\kappa,i\right) =
\left(T_{n},z\right)$. For the other degenerate states, $\left(T_{n},x\right)$
and $\left(T_{n},y\right)$, we can once again use the appropriate cyclic
transformations to swap valley axes.

% ------------------------------------------------------------------------------

We will summarize the findings of the preceding chapter here. To that end, we
will articulate a condensed form to collate all information required to specify
a particular state within a wider basis. We apply the following notational
nicety; by writing the symbol $\mathtt{I}$ in a fixed-width typewrite font, we
are denoting a single arbitrary basis state term of our full multi-valley
wavefunction. In particular, we imagine that $\mathtt{I}$ may function as a
label for the particular state, and as an subscript on a quantity denoting
membership in the collection of properties associated with the $\mathtt{I}$
state. We expect $\mathtt{I}$ to include all primary quantities necessary to
define a particular basis state. Information we might expect to be contained in
this structure would include the $\left(\kappa,i\right)$ (symmetry and
degeneracy), and $\varrho$ (single-axis function form) indices. It would also
necessarily include the full or implicit definitions of the naive envelope basis
functions in the form of discrete indices and continuous parameters. This
definition consists of the collection of discrete free parameters such as the
$m$, $l$, and $k$ indices, the $\eta$ and $\lambda$ continuous free parameters,
and the derived $\zeta$ acyclic transformation index. For convenience we can
imagine that, in addition to the necessary primary independent parameters, the
collection will also include derived secondary quantities used to construct the
full wavefunction, such the $\Omega$ index used in the axial Bloch functions. 

It should be made clear that the exact composition of elements in the collection
$\mathtt{I}$ is not static insofar as the number and nature of the values may
vary between different possible basis states we wish to explore. The structure
of $\mathtt{I}$ and will depend upon the type of symmetry denoted by $\kappa$
and the envelope function form denoted by $\varrho$. For example, all functions
include the $k$, $l$, and $m$ free discrete indices associated with the naive
envelope functions. For most states, $0\le m\le l$, and $0\le k$ are indices and
may take on arbitrary integer values satisfying the given constraints. However,
states using the $\parallel_{0}$ envelope function form do not include states
with odd $m$. Comparatively, states of the $\kappa = T_{n}$ symmetry using the
$\perp_{0}$ envelope function form only include states with $m$ odd but also
include the free index $h = \pm 1$ in addition to $\left(mlk\right)$. Similarly,
states with the $\kappa = E$ type of symmetry also require four independent
indices, where $\sigma = h$ than having its value determined by the symmetry
index $\kappa$ as it would be for the $\kappa = A_{n},T_{n}$ states.

We introduce this construct $\mathtt{I}$ as useful notation for denoting the
source of information and properties used within the analysis and software of
Cantankerous Impurity. For example, in this document it will see use in formulae
to compute the matrix elements of the Hamiltonian and overlap matrices. However,
we feel the primary value of introducing $\mathtt{I}$ lies in the manner in
which it will help guide and structure later software design. More precise and
rigorous definition of $\mathtt{I}$, if required, can be found in the relevant
software documentation.

We can articulate the following form for the multi-valley basis functions as a
combination of the so-called single-axis basis functions defined earlier in this
document.
\begin{equation}
\Psi_{\mathtt{I}}
    =
    \sum_{p}
        \left(\bm{\alpha}_{\mathtt{I}}\right)_{p}
        \psi_{\mathtt{I}}^{\left(p\right)}
,
\end{equation}
where $\alpha$ above has the same values as defined in subsection
\ref{subsection:symm_silicon}. We identify the set of indices
$p=\left(0,1,2\right)$ as given in that section with the valley-indices as
defined above going as $p=\left(z,x,y\right)$ respectively. Briefly, we may
reiterate that
\begin{equation}
\bm{\alpha}_{A_{n}}
    =
        \frac{1}{3}\left(1,1,1\right)
,
\end{equation}
as well as
\begin{align}
\bm{\alpha}_{\left(E,0\right)}
    &=
        \frac{1}{3}\left(1,\tau,\tau^{2}\right)
,&
\bm{\alpha}_{\left(E,1\right)}
    &=
        \frac{1}{3}\left(1,\tau^{2},\tau\right)
,
\end{align}
and
\begin{equation}
\bm{\alpha}_{\left(T_{n},i\right)}
    =
        \left(\delta_{ix},\delta_{iy},\delta_{iz}\right)
.
\end{equation}
As articulated previously for all multi-valley states there are multiple
independent constructions possible. Following the convention established earlier
in the document, we denote which construction is used with the state form index
$\varrho$. For most states, there are only two constructions of the single
valley basis possible, which we denote with $\varrho = \parallel_{0}$, or
$\varrho = \parallel_{1}$. For $\kappa = T_{n}$ there are five such
constructions possible. We denote the additional form indices using $\perp_{0}$,
$\perp_{1}$, and $\perp_{2}$. We refer to the $\parallel$ forms as
\emph{longitudinal} and the $\perp$ forms as \emph{transverse}, but the
distinction only exists and is only relevant for $T_{n}$ states.

%-------------------------------------------------------------------------------

% Another way we can articulate the multi-valley wavefunction is directly as the
% summation 
% \begin{equation}
% \Psi_{\mathtt{I}}
%     =
%     \sum_{q}
%         \left(\bm{\alpha}_{\mathtt{I}}\right)_{q}
%         \Phi_{\mathtt{I}}^{\left(p\right)}
%         \Theta_{ \Omega_\mathtt{I} }^{\left(q\right)}
% .
% \end{equation}

% \begin{equation}
% \Phi_{\mathtt{I}}^{\left(p\right)}
%     =
%     \sum_{p}
%         \left(\bm{\beta}_{\mathtt{I}}\right)_{pq}
%         \phi_{ 
%             \left(mlk\right)_{\mathtt{I}}
%             \left(\zeta_{\mathtt{I}}\right)_{pq} 
%         }^{\left(p\right)}
% \end{equation}



% The relevant values of $\bm{\beta}_{\mathtt{I}}$ and $\zeta_{\mathtt{I}}$



% \begin{equation}
% \xi
%     =
%         \sigma\left(-1\right)^{l+m}\im^{+m}
% \end{equation}


% where
% \begin{align}
% \bm{\beta}_{A_{n},\parallel_{0}}
% \bm{\beta}_{\left(E,0\right),\parallel_{0}}
% \bm{\beta}_{\left(E,1\right),\parallel_{0}}
% \bm{\beta}_{\left(T_{n},x\right),\parallel_{0}}
% &=
%     \hat{\bm{I}}_{3}
% ,
% \bm{\beta}_{A_{n},\parallel_{0}}
% \bm{\beta}_{\left(E,0\right),\parallel_{0}}
% \bm{\beta}_{\left(E,1\right),\parallel_{0}}
% &=
%     \hat{\bm{I}}_{3}
% ,
% \end{align}





% \begin{align}
% \bm{\beta}_{\left(T_{n},x\right),\parallel_{0}}
%     &=
%         \begin{pmatrix}
%             1 & 0 & 0 \\
%             0 & 0 & 0 \\
%             0 & 0 & 0
%         \end{pmatrix}
% ,
% \end{align}

% \begin{align}
% \bm{\beta}_{\left(T_{n},y\right),\parallel_{0}}
%     &=
%         \begin{pmatrix}
%             0 & 0 & 0 \\
%             0 & 1 & 0 \\
%             0 & 0 & 0
%         \end{pmatrix}
% ,
% \end{align}


% \begin{align}
% \bm{\beta}_{\left(T_{n},z\right),\parallel_{0}}
%     &=
%         \begin{pmatrix}
%             0 & 0 & 0 \\
%             0 & 0 & 0 \\
%             0 & 0 & 1
%         \end{pmatrix}
% ,
% \end{align}


% \begin{align}
% \zeta_{\parallel_{0}}
%     &=
%         \xi\left(-1\right)^{l+m}
% ,&
% \Omega_{\parallel_{0}}
%     &=
%         \omega\xi\left(-1\right)^{m}
% ,
% \end{align}




% \pagebreak






% \begin{align}
% \bm{\beta}_{A_{n},\parallel_{1}}
% \bm{\beta}_{\left(E,0\right),\parallel_{1}}
% \bm{\beta}_{\left(E,0\right),\parallel_{1}}
%     &=
%         \frac{1}{3}
%         \begin{pmatrix}
%             0   & \xi & 1   \\
%             1   & 0   & \xi \\
%             \xi & 1   & 0
%         \end{pmatrix}
% ,
% \end{align}



% \begin{align}
% \bm{\beta}_{\left(T_{n},x\right),\parallel_{1}}
%     &=
%         \begin{pmatrix}
%             0   & 0 & 0 \\
%             1   & 0 & 0 \\
%             \xi & 0 & 0
%         \end{pmatrix}
% ,
% \end{align}

% \begin{align}
% \bm{\beta}_{\left(T_{n},y\right),\parallel_{1}}
%     &=
%         \begin{pmatrix}
%             0 & \xi & 0 \\
%             0 & 0   & 0 \\
%             0 & 1   & 0
%         \end{pmatrix}
% ,
% \end{align}


% \begin{align}
% \bm{\beta}_{\left(T_{n},z\right),\parallel_{1}}
%     &=
%         \begin{pmatrix}
%             0 & 0 & 1   \\
%             0 & 0 & \xi \\
%             0 & 0 & 0
%         \end{pmatrix}
% ,
% \end{align}

% \begin{align}
% \left(\zeta_{\parallel_{1}}\right)_{p\left(p+1\right)}
%     &=
%         \left(-1\right)^{l}
% ,&
% \left(\zeta_{\parallel_{1}}\right)_{p\left(p+2\right)}
%     &=
%     \left(-1\right)^{l+m}
% ,&
% \Omega_{\parallel_{1}}
%     &=
%         \omega\left(-1\right)^{m}
% ,
% \end{align}





% \pagebreak





% \begin{align}
% \left(\zeta_{\perp_{0}}\right)_{p\left(p+1\right)}
%     &=
%         h
% ,&
% \left(\zeta_{\perp_{0}}\right)_{p\left(p+2\right)}
%     &=
%         -h
% ,&
% \Omega_{\perp_{0}}
%     &=
%         h\left(-1\right)^{m}
% ,
% \end{align}
    



% \begin{align}
% \left(\zeta_{\perp_{1}}\right)_{p\left(p+1\right)}
%         &=
%             -\left(-1\right)^{l}
%     ,&
% \left(\zeta_{\perp_{1}}\right)_{p\left(p+2\right)}
%         &=
%             -\left(-1\right)^{l+m}
% \end{align}


% \begin{align}
% \left(\zeta_{\perp_{2}}\right)_{p\left(p+1\right)}
%     &=
%         -\left(-1\right)^{l+m}
%     ,&
% \left(\zeta_{\perp_{2}}\right)_{p\left(p+2\right)}
% &=
%     -\left(-1\right)^{l}
% \end{align}




    










% \begin{align}
%     \bm{\beta}_{\left(T_{n},z\right),\perp_{0}}
%         &=
%             \begin{pmatrix}
%                 1 & 0 & 0 \\
%                 0 & -\zeta\xi\left(-1\right)^{l+m} & 0 \\
%                 0 & 0 & 0
%             \end{pmatrix}
%     ,&
%     \bm{\beta}_{\left(T_{n},z\right),\perp_{1}}
%         &=
%             \begin{pmatrix}
%                 0 & -\xi & 0   \\
%                 1 & 0 & 0 \\
%                 0 & 0 & 0
%             \end{pmatrix}
%     ,&
%     \bm{\beta}_{\left(T_{n},z\right),\perp_{2}}
%         &=
%             \begin{pmatrix}
%                 0 & 0 & 0   \\
%                 0 & 0 & 0 \\
%                 1 & -\xi\left(-1\right)^{m} & 0
%             \end{pmatrix}
%     ,
% \end{align}


% \begin{align}
%     \bm{\beta}_{\left(T_{n},x\right),\perp_{0}}
%         &=
%             \begin{pmatrix}
%                 0 & 0 & 0 \\
%                 0 & 1 & 0 \\
%                 0 & 0 & -\zeta\xi\left(-1\right)^{l+m} \\
%             \end{pmatrix}
%     ,&
%     \bm{\beta}_{\left(T_{n},x\right),\perp_{1}}
%         &=
%             \begin{pmatrix}
%                 0 & 0 & 0   \\
%                 0 & 0 & -\xi \\
%                 0 & 1 & 0
%             \end{pmatrix}
%     ,&
%     \bm{\beta}_{\left(T_{n},x\right),\perp_{2}}
%         &=
%             \begin{pmatrix}
%                 0 & 1 & -\xi\left(-1\right)^{m} \\
%                 0 & 0 & 0                       \\
%                 0 & 0 & 0
%             \end{pmatrix}
%     ,
% \end{align}


% \begin{align}
%     \bm{\beta}_{\left(T_{n},y\right),\perp_{0}}
%         &=
%             \begin{pmatrix}
%                 -\zeta\xi\left(-1\right)^{l+m} & 0 & 0 \\
%                 0                          & 0 & 0 \\
%                 0                          & 0 & 1 \\
%             \end{pmatrix}
%     ,&
%     \bm{\beta}_{\left(T_{n},y\right),\perp_{1}}
%         &=
%             \begin{pmatrix}
%                 0    & 0 & 1 \\
%                 0    & 0 & 0 \\
%                 -\xi & 0 & 0
%             \end{pmatrix}
%     ,&
%     \bm{\beta}_{\left(T_{n},y\right),\perp_{2}}
%         &=
%             \begin{pmatrix}
%                 0                       & 0 & 0 \\
%                 -\xi\left(-1\right)^{m} & 0 & 1 \\
%                 0                       & 0 & 0
%             \end{pmatrix}
%     ,
% \end{align}













Another way we can articulate the multi-valley wavefunction is directly as the
summation 
% \begin{equation}
% \Psi_{\mathtt{I}}
%     =
%     \sum_{pq\in\left(xyz\right)}
%         \left(\bm{\beta}_{\mathtt{I}}\right)_{pq}
%         \phi_{ 
%             \left(mlk\right)_{\mathtt{I}}
%             \left(\zeta_{\mathtt{I}}\right)_{pq} 
%         }^{\left(p\right)}
%         \Theta_{ \Omega_\mathtt{I} }^{\left(q\right)}
% .
% \end{equation}
\begin{equation}
\Psi_{\mathtt{I}}
    =
        \sum_{q}
            \Phi_{\mathtt{I}}^{\left(q\right)}
            \Theta_{ \Omega_\mathtt{I} }^{\left(q\right)}
.
\end{equation}
% The relevant values of $\bm{\beta}_{\mathtt{I}}$ and $\zeta_{\mathtt{I}}$


\begin{equation}
\Phi_{\mathtt{I}}^{\left(q\right)}
    =
        \sum_{p}
            \left(\bm{\beta}_{\mathtt{I}}\right)_{pq}
            \phi_{ 
                \left(mlk\right)_{\mathtt{I}}
                \left(\zeta_{\mathtt{I}}\right)_{pq} 
            }^{\left(p\right)}
.
\end{equation}



\begin{equation}
\xi
    =
        \sigma\left(-1\right)^{l+m}\im^{+m}
\end{equation}


where
\begin{align}
\left(\bm{\beta}_{\left(\kappa,i\right),\parallel_{0}}\right)_{pq}
    &=
        \left(\bm{\alpha}_{\left(\kappa,i\right)}\right)_{q}
        \delta_{pq}
,&
\zeta_{\parallel_{0}}
    &=
        \xi\left(-1\right)^{l+m}
,&
\Omega_{\parallel_{0}}
    &=
        \omega\xi\left(-1\right)^{m}
,
\end{align}




\pagebreak






\begin{align}
\bm{\beta}_{A_{n},\parallel_{1}}
    &=
        \frac{1}{3}
        \begin{pmatrix}
            0   & \xi & 1   \\
            1   & 0   & \xi \\
            \xi & 1   & 0
        \end{pmatrix}
,
\end{align}

\begin{align}
\bm{\beta}_{\left(E,0\right),\parallel_{1}}
    &=
        \frac{1}{3}
        \begin{pmatrix}
            0       & \tau^{2}\xi & 1   \\
            \tau    & 0           & \xi \\
            \tau\xi & \tau^{2}    & 0
        \end{pmatrix}
,
\end{align}

\begin{align}
\bm{\beta}_{\left(E,0\right),\parallel_{1}}
    &=
        \frac{1}{3}
        \begin{pmatrix}
            0           & \tau\xi & 1   \\
            \tau^{2}    & 0       & \xi \\
            \tau^{2}\xi & \tau    & 0
        \end{pmatrix}
,
\end{align}



\begin{align}
\bm{\beta}_{\left(T_{n},x\right),\parallel_{1}}
    &=
        \begin{pmatrix}
            0   & 0 & 0 \\
            1   & 0 & 0 \\
            \xi & 0 & 0
        \end{pmatrix}
,
\end{align}

\begin{align}
\bm{\beta}_{\left(T_{n},y\right),\parallel_{1}}
    &=
        \begin{pmatrix}
            0 & \xi & 0 \\
            0 & 0   & 0 \\
            0 & 1   & 0
        \end{pmatrix}
,
\end{align}


\begin{align}
\bm{\beta}_{\left(T_{n},z\right),\parallel_{1}}
    &=
        \begin{pmatrix}
            0 & 0 & 1   \\
            0 & 0 & \xi \\
            0 & 0 & 0
        \end{pmatrix}
,
\end{align}

\begin{align}
\left(\zeta_{\parallel_{1}}\right)_{\left(q+1\right)q}
    &=
        \left(-1\right)^{l}
,&
\left(\zeta_{\parallel_{1}}\right)_{\left(q+2\right)q}
    &=
        \left(-1\right)^{m}
        \left(\zeta_{\parallel_{1}}\right)_{q\left(q+1\right)}
,&
\Omega_{\parallel_{1}}
    &=
        \omega\left(-1\right)^{m}
,
\end{align}





\pagebreak





\begin{align}
\left(\zeta_{\perp_{0}}\right)_{\left(q+1\right)\left(q+1\right)}
    &=
        h
,&
\left(\zeta_{\perp_{0}}\right)_{\left(q+2\right)\left(q+2\right)}
    &=
        \left(-1\right)^{m}
        \left(\zeta_{\perp_{0}}\right)_{\left(q+1\right)\left(q+1\right)}
,&
\Omega_{\perp_{0}}
    &=
        h\left(-1\right)^{m}
,
\end{align}
    



\begin{align}
\left(\zeta_{\perp_{1}}\right)_{\left(q+1\right)\left(q+2\right)}
        &=
            -\left(-1\right)^{l}
    ,&
\left(\zeta_{\perp_{1}}\right)_{\left(q+2\right)\left(q+1\right)}
        &=
            \left(-1\right)^{m}
            \left(\zeta_{\perp_{1}}\right)_{\left(q+1\right)\left(q+2\right)}
\end{align}


\begin{align}
\left(\zeta_{\perp_{2}}\right)_{q\left(q+1\right)}
    &=
        -\left(-1\right)^{l+m}
,&
\left(\zeta_{\perp_{2}}\right)_{q\left(q+2\right)}
&=
    \left(-1\right)^{m}
    \left(\zeta_{\perp_{2}}\right)_{q\left(q+1\right)}
\end{align}




    










\begin{align}
    \bm{\beta}_{\left(T_{n},z\right),\perp_{0}}
        &=
            \begin{pmatrix}
                1 & 0 & 0 \\
                0 & -\zeta\xi\left(-1\right)^{l+m} & 0 \\
                0 & 0 & 0
            \end{pmatrix}
    ,&
    \bm{\beta}_{\left(T_{n},z\right),\perp_{1}}
        &=
            \begin{pmatrix}
                0 & -\xi & 0   \\
                1 & 0 & 0 \\
                0 & 0 & 0
            \end{pmatrix}
    ,&
    \bm{\beta}_{\left(T_{n},z\right),\perp_{2}}
        &=
            \begin{pmatrix}
                0 & 0 & 0   \\
                0 & 0 & 0 \\
                1 & -\xi\left(-1\right)^{m} & 0
            \end{pmatrix}
    ,
\end{align}


\begin{align}
    \bm{\beta}_{\left(T_{n},x\right),\perp_{0}}
        &=
            \begin{pmatrix}
                0 & 0 & 0 \\
                0 & 1 & 0 \\
                0 & 0 & -\zeta\xi\left(-1\right)^{l+m} \\
            \end{pmatrix}
    ,&
    \bm{\beta}_{\left(T_{n},x\right),\perp_{1}}
        &=
            \begin{pmatrix}
                0 & 0 & 0   \\
                0 & 0 & -\xi \\
                0 & 1 & 0
            \end{pmatrix}
    ,&
    \bm{\beta}_{\left(T_{n},x\right),\perp_{2}}
        &=
            \begin{pmatrix}
                0 & 1 & -\xi\left(-1\right)^{m} \\
                0 & 0 & 0                       \\
                0 & 0 & 0
            \end{pmatrix}
    ,
\end{align}


\begin{align}
    \bm{\beta}_{\left(T_{n},y\right),\perp_{0}}
        &=
            \begin{pmatrix}
                -\zeta\xi\left(-1\right)^{l+m} & 0 & 0 \\
                0                          & 0 & 0 \\
                0                          & 0 & 1 \\
            \end{pmatrix}
    ,&
    \bm{\beta}_{\left(T_{n},y\right),\perp_{1}}
        &=
            \begin{pmatrix}
                0    & 0 & 1 \\
                0    & 0 & 0 \\
                -\xi & 0 & 0
            \end{pmatrix}
    ,&
    \bm{\beta}_{\left(T_{n},y\right),\perp_{2}}
        &=
            \begin{pmatrix}
                0                       & 0 & 0 \\
                -\xi\left(-1\right)^{m} & 0 & 1 \\
                0                       & 0 & 0
            \end{pmatrix}
    ,
\end{align}









In short, we imagine 




we can imagine the collection $\mathtt{I}$ includes at least 3, and
sometimes 4, discrete free parameters. We can denote those parameters by
$\left(mlkh\right)$ with h representing the value of the fourth discrete free
parameters if present and ignored otherwise.To fully define the complete
multi-valley 




\begin{align}
\Psi_{\varrho,mlk}^{\left(A_{n}\right)}
    &=
        \frac{1}{3}
        \psi_{\varrho,mlk\sigma^{\left(A_{n}\right)}\left(+1\right)}
        ^{\left(z\right)}
            +
        \frac{1}{3}
        \psi_{\varrho,mlk\sigma^{\left(A_{n}\right)}\left(+1\right)}
        ^{\left(x\right)}
            +
        \frac{1}{3}
        \psi_{\varrho,mlk\sigma^{\left(A_{n}\right)}\left(+1\right)}
        ^{\left(y\right)}
\label{wavefunction:A}
,\\
\Psi_{\varrho,mlk\sigma}^{\left(E,0\right)}
    &=
        \frac{1}{3}
        \psi_{\varrho,mlk\sigma\left(+1\right)}^{\left(z\right)}
            +
        \frac{\tau}{3}
        \psi_{\varrho,mlk\sigma\left(+1\right)}^{\left(x\right)}
            +
        \frac{\tau^{2}}{3}
        \psi_{\varrho,mlk\sigma\left(+1\right)}^{\left(y\right)}
\label{wavefunction:E:0}
,\\
\Psi_{\varrho,mlk\sigma}^{\left(E,1\right)}
    &=
        \frac{1}{3}
        \psi_{\varrho,mlk\sigma\left(+1\right)}^{\left(z\right)}
            +
        \frac{\tau^{2}}{3}
        \psi_{\varrho,mlk\sigma\left(+1\right)}^{\left(x\right)}
            +
        \frac{\tau}{3}
        \psi_{\varrho,mlk\sigma\left(+1\right)}^{\left(y\right)}
\label{wavefunction:E:1}
.\\
\Psi_{\varrho,mlk}^{\left(T_{n},q\right)}
    &=
        \delta_{qz}
        \psi_{\varrho,mlk\sigma^{\left(T_{n}\right)}\left(-1\right)}
        ^{\left(z\right)}
            +
        \delta_{qx}
        \psi_{\varrho,mlk\sigma^{\left(T_{n}\right)}\left(-1\right)}
        ^{\left(x\right)}
            +
        \delta_{qy}
        \psi_{\varrho,mlk\sigma^{\left(T_{n}\right)}\left(-1\right)}
        ^{\left(y\right)}
\label{wavefunction:T:q}
.
\end{align}
\begin{align}
\psi_{\parallel_{0},m_{e}lk\sigma\omega}^{\left(z\right)}
        & \coloneqq
            \phi_{klm_{e}\zeta}^{(z)}
            \Theta^{\left(z\right)}
            _{\sigma\omega\left(-1\right)^{l+m_{e}/2}}
\tag{\ref{envelope:longitudinal:z:0}+\ref{single-axis:longitudinal:z:0}}
,\\
\psi_{\parallel_{1},mlk\sigma\omega}^{\left(z\right)}
    & \coloneqq
        \left\{
            \phi_{klm\zeta_{0}}^{(x)}
            +
            \sigma\left(-1\right)^{l+m}\im^{+m}
            \phi_{klm\zeta_{1}}^{(y)}
        \right\}
        \Theta^{\left(z\right)}
        _{\omega\left(-1\right)^{m}}
\tag{\ref{envelope:longitudinal:z:1}+\ref{single-axis:longitudinal:z:1}}
.
\end{align}


\begin{align}
\Psi_{\perp_{0},m_{o}lk\zeta}^{\left(T_{n},z\right)}
    &=
        \phi_{m_{o}lk\zeta}^{\left(x\right)}
        \Theta^{\left(x\right)}_{\zeta\left(-1\right)^{l}}
        -
        \zeta\sigma^{\left(T_{n}\right)}\im^{+m}
        \phi_{m_{o}lk\bar{\zeta}}^{\left(y\right)}
        \Theta^{\left(y\right)}_{\zeta\left(-1\right)^{l}}        
,\\
\Psi_{\perp_{1},mlk}^{\left(T_{n},z\right)}
    &=
        % \phi_{mlk\left(-1\right)^{l+1}}^{\left(y\right)}
        \phi_{mlk\bar{\zeta}_{0}}^{\left(y\right)}
        \Theta^{\left(x\right)}_{\left(-1\right)^{m+1}}
        -
        \sigma^{\left(T_{n}\right)}\left(-1\right)^{l+m}\im^{+m}
        % \phi_{mlk\left(-1\right)^{l+m+1}}^{\left(x\right)}
        \phi_{mlk\bar{\zeta}_{1}}^{\left(x\right)}
        \Theta^{\left(y\right)}_{\left(-1\right)^{m+1}}        
,\\
\Psi_{\perp_{2},mlk}^{\left(T_{n},z\right)}
    &=
        \phi_{mlk\bar{\zeta}_{1}}^{\left(z\right)}
        \Theta^{\left(x\right)}_{\left(-1\right)^{m}}
        -
        \sigma^{\left(T_{n}\right)}\left(-1\right)^{l}\im^{+m}
        \phi_{mlk\bar{\zeta}_{0}}^{\left(z\right)}
        \Theta^{\left(y\right)}_{\left(-1\right)^{m}}
.
\end{align}

% \begin{align}
%     \Psi_{\perp_{0},m_{o}lk\zeta}^{\left(T_{n},z\right)}
%         &=
%             \Phi_{\perp_{0},m_{o}lk\zeta}^{\left(x\right)}
%             \Theta^{\left(x\right)}_{\Omega}
%             -
%             \zeta\sigma^{\left(T_{n}\right)}\im^{+m}
%             \Phi_{\perp_{0},m_{o}lk\zeta}^{\left(y\right)}
%             \Theta^{\left(y\right)}_{\Omega}
%     ,&
%     \Omega
%         &=
%             \zeta\left(-1\right)^{l}
%     ,\\
%     \Psi_{\perp_{1},mlk}^{\left(T_{n},z\right)}
%         &=
%             % \phi_{mlk\left(-1\right)^{l+1}}^{\left(y\right)}
%             \Phi_{\perp_{1},mlk}^{\left(x\right)}
%             \Theta^{\left(x\right)}_{\Omega}
%             -
%             \sigma^{\left(T_{n}\right)}\left(-1\right)^{l}\im^{-m}
%             % \phi_{mlk\left(-1\right)^{l+m+1}}^{\left(x\right)}
%             \Phi_{\perp_{1},mlk}^{\left(y\right)}
%             \Theta^{\left(y\right)}_{\Omega}
%     ,&
%     \Omega
%         &=
%             \left(-1\right)^{m+1}
%     ,\\
%     \Psi_{\perp_{2},mlk}^{\left(T_{n},z\right)}
%         &=
%             \Phi_{\perp_{2},mlk}^{\left(x\right)}
%             \Theta^{\left(x\right)}_{\Omega}
%             -
%             \sigma^{\left(T_{n}\right)}\left(-1\right)^{l}\im^{+m}
%             \Phi_{\perp_{2},mlk}^{\left(u\right)}
%             \Theta^{\left(y\right)}_{\Omega}
%     &
%     \Omega
%         &=
%             \left(-1\right)^{m}
%     .
% \end{align}


% \begin{equation}
% \psi_{\mathtt{I}}^{\left(p\right)}
%     =
%     \sum_{qr}
%         \beta_{qr}^{\left(\varrho,p\right)}
%         \phi_{\mathtt{I}}^{\left(q\right)}
%         \Theta_{\mathtt{I}}^{\left(r\right)}
% .
% \end{equation}


% We define the following as our \emph{naive envelope basis function} with
% anisotropy included as follows,
% \begin{equation}
%   \Phi_{klm;\zeta;\eta\lambda;\bm{r}}^{(Wq)}
%   \coloneqq
%   F_{lk;\eta\lambda;r\chi_{q}}^{\left(W\right)}
%   P_{lm;\lambda;\chi_{q}}
%   Q_{m\zeta;\phi_{q}}
%   ,
% \tag{\ref{env:naive}}
% \end{equation}


% ------------------------------------------------------------------------------

\begin{align}
\mathcal{H}
    &=
        \int \diff \bm{r}\,
        \left(\Psi_{\mathtt{L}}\right)^{\dagger}
        \hat{U}_{\bm{r}}\,
        \Psi_{\mathtt{R}}
,\\
    &=
        \int \diff \bm{r}\,
        \left\{
            \sum_{q_{\mathtt{L}}}
            \Phi_{\mathtt{L}}^{\left(q_{\mathtt{L}}\right)}
            \Theta_{ \Omega_\mathtt{L} }^{\left(q_{\mathtt{L}}\right)}
        \right\}^{\dagger}
        \hat{U}_{\bm{r}}\,
        \left\{
            \sum_{q_{\mathtt{R}}}
            \Phi_{\mathtt{R}}^{\left(q_{\mathtt{R}}\right)}
            \Theta_{ \Omega_\mathtt{R} }^{\left(q_{\mathtt{R}}\right)}
        \right\}
,\\
    &=
        \int \diff \bm{r}\,
        \sum_{q_{\mathtt{L}}}
        \sum_{q_{\mathtt{R}}}
            \tilde{\Phi_{\mathtt{L}}}^{\left(q_{\mathtt{L}}\right)}
        \hat{U}_{\bm{r}}\,
        \Phi_{\mathtt{R}}^{\left(q_{\mathtt{R}}\right)}
        \Theta_{ \Omega_\mathtt{L} }^{\left(-q_{\mathtt{L}}\right)}
        \Theta_{ \Omega_\mathtt{R} }^{\left(+q_{\mathtt{R}}\right)}
,\\
    &=
        \mathcal{H}_{\parallel}
            +
        \mathcal{H}_{\perp}
.
\end{align}

\begin{align}
\mathcal{H}_{\parallel}
    &=
        \int \diff \bm{r}\,
        \sum_{q=x,y,z}
            \Phi_{\mathtt{L}}^{\left(q\right)*}
            \hat{U}_{\bm{r}}\,
            \Phi_{\mathtt{R}}^{\left(q\right)}
            \Theta_{ \Omega_\mathtt{L} }^{\left(-q\right)}
            \Theta_{ \Omega_\mathtt{R} }^{\left(q\right)}
,\\
    &=
        \int \diff \bm{r}\,
        \varphi_{\parallel}
        \Theta_{ \Omega_\mathtt{L} }^{\left(-z\right)}
        \Theta_{ \Omega_\mathtt{R} }^{\left(z\right)}
.
\end{align}
\begin{align}
\varphi_{\parallel}
    &=
        \Phi_{\mathtt{L}}^{\left(z\right)*}
        \hat{U}_{\bm{r}}\,
        \Phi_{\mathtt{R}}^{\left(z\right)}
            +
        \left(zxy\right)
        \Phi_{\mathtt{L}}^{\left(x\right)*}
        \hat{U}_{\bm{r}}\,
        \Phi_{\mathtt{R}}^{\left(x\right)}
            +
        \left(yzx\right)
        \Phi_{\mathtt{L}}^{\left(y\right)*}
        \hat{U}_{\bm{r}}\,
        \Phi_{\mathtt{R}}^{\left(y\right)}
,\\
    &=
        \sum_{p_{\mathtt{L}},p_{\mathtt{R}}}
            \left(
                \bm{\beta}^{\left(z\right)}_{p_{\mathtt{L}}p_{\mathtt{R}}}
                \hat{U}_{xyz}
                \phi_{zz}^{\left(p_{\mathtt{L}}p_{\mathtt{R}}\right)}
                    +
                \bm{\beta}^{\left(x\right)}_{p_{\mathtt{L}}p_{\mathtt{R}}}
                \hat{U}_{zxy}
                \phi_{xx}^{\left(p_{\mathtt{L}}p_{\mathtt{R}}\right)}
                    +
                \bm{\beta}^{\left(y\right)}_{p_{\mathtt{L}}p_{\mathtt{R}}}
                \hat{U}_{yzx}
                \phi_{yy}^{\left(p_{\mathtt{L}}p_{\mathtt{R}}\right)}
            \right)
.
\end{align}

\begin{align}
\bm{\beta}_{p_{\mathtt{L}}p_{\mathtt{R}}}
^{\left(z-\Delta\right)}
    &=
        \left(\bm{\beta}_{\mathtt{L}}\right)
        _{\left(p_{\mathtt{L}}-\Delta\right)\left(z-\Delta\right)}^{*}
        \left(\bm{\beta}_{\mathtt{R}}\right)
        _{\left(p_{\mathtt{R}}-\Delta\right)\left(z-\Delta\right)}
,\\
\phi_{\left(z-\Delta\right)\left(z-\Delta\right)}
^{\left(p_{\mathtt{L}}p_{\mathtt{R}}\right)}
    &=
        \tilde{\phi}_{
            \mathtt{L},
            \left(\zeta_\mathtt{L}\right)
            _{\left(p_{\mathtt{L}}-\Delta\right)\left(z-\Delta\right)}
        }^{\left(p_{\mathtt{L}}\right)}
        \phi_{
            \mathtt{R},
            \left(\zeta_\mathtt{R}\right)
            _{\left(p_{\mathtt{R}}-\Delta\right)\left(z-\Delta\right)}
        }^{\left(p_{\mathtt{R}}\right)}
,\\
\zeta_{\left(z-\Delta\right)\left(z-\Delta\right)}
^{\left(p_{\mathtt{L}}p_{\mathtt{R}}\right)}
    &=
        \left(\zeta_\mathtt{L}\right)
        _{\left(p_{\mathtt{L}}-\Delta\right)\left(z-\Delta\right)}^{*}
        \left(\zeta_\mathtt{R}\right)
        _{\left(p_{\mathtt{R}}-\Delta\right)\left(z-\Delta\right)}
.
\end{align}

\begin{align}
% \varphi_{\parallel}
%     &=
%         \varphi_{\parallel,zz}
%             +
%         \varphi_{\parallel,xx}
%             +
%         \varphi_{\parallel,yy}
%             +
%         \sum_{p_{\mathtt{L}},p_{\mathtt{R}}}
%         \left(
%             \left(\bm{\beta}_{\mathtt{L}\mathtt{R}}\right)^{\left(z\right)}
%             _{p_{\mathtt{L}}p_{\mathtt{R}}}
%             \hat{U}_{xyz}
%                 +
%             \left(\bm{\beta}_{\mathtt{L}\mathtt{R}}\right)^{\left(x\right)}
%             _{p_{\mathtt{L}}p_{\mathtt{R}}}
%             \hat{U}_{zxy}
%                 +
%             \left(\bm{\beta}_{\mathtt{L}\mathtt{R}}\right)
%             _{p_{\mathtt{L}}p_{\mathtt{R}}}
%             ^{\left(y\right)}
%             \hat{U}_{yzx}
%         \right)
%         \tilde{\phi}_{\mathtt{L}}^{\left(p_{\mathtt{L}}\right)}
%         \phi_{\mathtt{R}}^{\left(p_{\mathtt{R}}\right)}
% ,\\
\varphi_{\parallel,zz}
    &=
        \bm{\beta}^{\left(z\right)}_{zz}
        \hat{U}_{xyz}
        \phi_{zz}^{\left(zz\right)}
            +
        \bm{\beta}^{\left(x\right)}_{zz}
        \hat{U}_{zxy}
        \phi_{xx}^{\left(zz\right)}
            +
        \bm{\beta}^{\left(y\right)}_{zz}
        \hat{U}_{yzx}
        \phi_{yy}^{\left(zz\right)}
,\\
\varphi_{\parallel,xx}
    &=
        \bm{\beta}^{\left(z\right)}_{xx}
        \hat{U}_{xyz}
        \phi_{zz}^{\left(xx\right)}
            +
        \bm{\beta}^{\left(x\right)}_{xx}
        \hat{U}_{zxy}
        \phi_{xx}^{\left(xx\right)}
            +
        \bm{\beta}^{\left(y\right)}_{xx}
        \hat{U}_{yzx}
        \phi_{yy}^{\left(xx\right)}
,\\
\varphi_{\parallel,yy}
    &=
        \bm{\beta}^{\left(z\right)}_{yy}
        \phi_{zz}^{\left(yy\right)}
        \hat{U}_{xyz}
            +
        \bm{\beta}^{\left(x\right)}_{yy}
        \hat{U}_{zxy}
        \phi_{xx}^{\left(yy\right)}
            +
        \bm{\beta}^{\left(y\right)}_{yy}
        \hat{U}_{yzx}
        \phi_{yy}^{\left(yy\right)}
\end{align}









    

\begin{align}
\left(yxz\right)\varphi_{\parallel,zz}
    &=
        \bm{\beta}^{\left(z\right)}_{zz}
        \left(\zeta_{zz}^{\left(zz\right)}\hat{U}_{yxz}\right)
        \phi_{zz}^{\left(zz\right)}
        +
        \left(
            \bm{\beta}^{\left(y\right)}_{zz}
            \zeta_{yy}^{\left(zz\right)}
            \hat{U}_{xzy}
        \right)
        \phi_{xx}^{\left(zz\right)}
            +
        \left(
            \bm{\beta}^{\left(x\right)}_{zz}
            \zeta_{xx}^{\left(zz\right)}
            \hat{U}_{zyx}
        \right)
        \phi_{yy}^{\left(zz\right)}
,\\
\left(yxz\right)\varphi_{\parallel,xx}
    &=
        \bm{\beta}^{\left(z\right)}_{xx}
        \hat{U}_{yxz}
        \phi_{zz}^{\left(xx\right)}
            +
        \bm{\beta}^{\left(x\right)}_{xx}
        \hat{U}_{zyx}
        \phi_{xx}^{\left(xx\right)}
            +
        \bm{\beta}^{\left(y\right)}_{xx}
        \hat{U}_{xzy}
        \phi_{yy}^{\left(xx\right)}
,\\
\left(yxz\right)\varphi_{\parallel,yy}
    &=
        \bm{\beta}^{\left(z\right)}_{yy}
        \hat{U}_{yxz}
        \phi_{zz}^{\left(yy\right)}
            +
        \bm{\beta}^{\left(x\right)}_{yy}
        \hat{U}_{zyx}
        \phi_{xx}^{\left(yy\right)}
            +
        \bm{\beta}^{\left(y\right)}_{yy}
        \hat{U}_{xzy}
        \phi_{yy}^{\left(yy\right)}
\end{align}
    


\begin{align}
\left(\frac{\left(xyz\right)+\left(yxz\right)}{2}\right)\varphi_{\parallel,zz}
    &=
        \bm{\beta}^{\left(z\right)}_{zz}
        \left(
            \frac{\hat{U}_{xyz} + \zeta_{zz}^{\left(zz\right)}\hat{U}_{yxz}}{2}
        \right)
        \phi_{zz}^{\left(zz\right)}
        +
        \left(\frac{
                \bm{\beta}^{\left(x\right)}_{zz}
                \hat{U}_{zxy}
                    +
                \bm{\beta}^{\left(y\right)}_{zz}
                \zeta_{yy}^{\left(zz\right)}
                \hat{U}_{xzy}
        }{2}\right)
        \phi_{xx}^{\left(zz\right)}
            +
        \left(\frac{
            \bm{\beta}^{\left(y\right)}_{zz}
            \hat{U}_{yzx}
                +
            \bm{\beta}^{\left(x\right)}_{zz}
            \zeta_{xx}^{\left(zz\right)}
            \hat{U}_{zyx}
        }{2}
        \right)
        \phi_{yy}^{\left(zz\right)}
% ,\\
\end{align}













\begin{align}
\mathcal{H}_{\perp}
    &=
        \int \diff \bm{r}\,
        \sum_{q=x,y,z}
        \left[
            \left\{
                \Phi_{\mathtt{L}}^{\left(q\right)}
            \right\}^{\dagger}
            \hat{U}_{\bm{r}}\,
            \Phi_{\mathtt{R}}^{\left(q+1\right)}
            \Theta_{ \Omega_\mathtt{L} }^{\left(-q\right)}
            \Theta_{ \Omega_\mathtt{R} }^{\left(q+1\right)}
                +
            \left\{
                \Phi_{\mathtt{L}}^{\left(q\right)}
            \right\}^{\dagger}
            \hat{U}_{\bm{r}}\,
            \Phi_{\mathtt{R}}^{\left(q-1\right)}
            \Theta_{ \Omega_\mathtt{L} }^{\left(-q\right)}
            \Theta_{ \Omega_\mathtt{R} }^{\left(q-1\right)}
        \right]
,\\
%     &=
%         \int \diff \bm{r}\,
%         \left[
%             \hat{U}_{\bm{r}}\,
%             \Phi_{\mathtt{L}}^{\left(x\right)*}
%             \Phi_{\mathtt{R}}^{\left(y\right)}
%                 +
%             \left(zxy\right)
%             \hat{U}_{\bm{r}}\,
%             \Phi_{\mathtt{L}}^{\left(y\right)*}
%             \Phi_{\mathtt{R}}^{\left(z\right)}
%                 +
%             \left(yzx\right)
%             \hat{U}_{\bm{r}}\,
%             \Phi_{\mathtt{L}}^{\left(z\right)*}
%             \Phi_{\mathtt{R}}^{\left(x\right)}
%         \right]
%         \Theta_{ \Omega_\mathtt{L} }^{\left(-x\right)}
%         \Theta_{ \Omega_\mathtt{R} }^{\left(y\right)}
%                 +
%         \left[
%             \hat{U}_{\bm{r}}\,
%             \Phi_{\mathtt{L}}^{\left(y\right)*}
%             \Phi_{\mathtt{R}}^{\left(x\right)}
%                 +
%             \left(zxy\right)
%             \hat{U}_{\bm{r}}\,
%             \Phi_{\mathtt{L}}^{\left(z\right)*}
%             \Phi_{\mathtt{R}}^{\left(y\right)}
%                 +
%             \left(yzx\right)
%             \hat{U}_{\bm{r}}\,
%             \Phi_{\mathtt{L}}^{\left(x\right)*}
%             \Phi_{\mathtt{R}}^{\left(z\right)}
%         \right]
%         \Theta_{ \Omega_\mathtt{L} }^{\left(-y\right)}
%         \Theta_{ \Omega_\mathtt{R} }^{\left(x\right)}
% ,\\
%     &=
%         \int \diff \bm{r}\,
%         \left\{
%             \left(yxz\right)\left[
%                 \hat{U}_{\bm{r}}\,
%                 \Phi_{\mathtt{L}}^{\left(x\right)*}
%                 \Phi_{\mathtt{R}}^{\left(y\right)}
%                     +
%                 \left(zxy\right)
%                 \hat{U}_{\bm{r}}\,
%                 \Phi_{\mathtt{L}}^{\left(y\right)*}
%                 \Phi_{\mathtt{R}}^{\left(z\right)}
%                     +
%                 \left(yzx\right)
%                 \hat{U}_{\bm{r}}\,
%                 \Phi_{\mathtt{L}}^{\left(z\right)*}
%                 \Phi_{\mathtt{R}}^{\left(x\right)}
%             \right]
%                 +
%             \left[
%                 \hat{U}_{\bm{r}}\,
%                 \Phi_{\mathtt{L}}^{\left(y\right)*}
%                 \Phi_{\mathtt{R}}^{\left(x\right)}
%                     +
%                 \left(zxy\right)
%                 \hat{U}_{\bm{r}}\,
%                 \Phi_{\mathtt{L}}^{\left(z\right)*}
%                 \Phi_{\mathtt{R}}^{\left(y\right)}
%                     +
%                 \left(yzx\right)
%                 \hat{U}_{\bm{r}}\,
%                 \Phi_{\mathtt{L}}^{\left(x\right)*}
%                 \Phi_{\mathtt{R}}^{\left(z\right)}
%             \right]
%         \right\}
%         \Theta_{ \Omega_\mathtt{L} }^{\left(-y\right)}
%         \Theta_{ \Omega_\mathtt{R} }^{\left(x\right)}
% ,\\
    &=
        \int \diff \bm{r}\,
        \left(\varphi_{\perp,0}+\left(yxz\right)\varphi_{\perp,0}\right)
        \Theta_{ \Omega_\mathtt{L} }^{\left(-y\right)}
        \Theta_{ \Omega_\mathtt{R} }^{\left(x\right)}
.
\end{align}
\begin{align}
\varphi_{\perp,0}
    &=
        \hat{U}_{\bm{r}}\,
        \Phi_{\mathtt{L}}^{\left(y\right)*}
        \Phi_{\mathtt{R}}^{\left(x\right)}
            +
        \left(zxy\right)
        \hat{U}_{\bm{r}}\,
        \Phi_{\mathtt{L}}^{\left(z\right)*}
        \Phi_{\mathtt{R}}^{\left(y\right)}
            +
        \left(yzx\right)
        \hat{U}_{\bm{r}}\,
        \Phi_{\mathtt{L}}^{\left(x\right)*}
        \Phi_{\mathtt{R}}^{\left(z\right)}
,\\
        % \left(yxz\right)
        % \left\{
        %     \hat{U}_{\bm{r}}
        %         +
        %     \left(zxy\right)
        %     \hat{U}_{\bm{r}}
        %         +
        % \right\}
        % \Phi_{\mathtt{L}}^{\left(x\right)*}
        % \Phi_{\mathtt{R}}^{\left(y\right)}        
        %     +
        % \left(zyx\right)
        % \hat{U}_{\bm{r}}\,
        % \Phi_{\mathtt{L}}^{\left(y\right)*}
        % \Phi_{\mathtt{R}}^{\left(z\right)}
        %     +
        % \left(xzy\right)
        % \hat{U}_{\bm{r}}\,
        % \Phi_{\mathtt{L}}^{\left(z\right)*}
        % \Phi_{\mathtt{R}}^{\left(x\right)}
        %     +
        % \hat{U}_{\bm{r}}\,
        % \Phi_{\mathtt{L}}^{\left(y\right)*}
        % \Phi_{\mathtt{R}}^{\left(x\right)}
        %     +
        % \left(zxy\right)
        % \hat{U}_{\bm{r}}\,
        % \Phi_{\mathtt{L}}^{\left(z\right)*}
        % \Phi_{\mathtt{R}}^{\left(y\right)}
        %     +
        % \left(yzx\right)
        % \hat{U}_{\bm{r}}\,
        % \Phi_{\mathtt{L}}^{\left(x\right)*}
        % \Phi_{\mathtt{R}}^{\left(z\right)}
\varphi_{\perp,1}
    &=
        \hat{U}_{\bm{r}}\,
        \Phi_{\mathtt{L}}^{\left(x\right)*}
        \Phi_{\mathtt{R}}^{\left(y\right)}
            +
        \left(zxy\right)
        \hat{U}_{\bm{r}}\,
        \Phi_{\mathtt{L}}^{\left(y\right)*}
        \Phi_{\mathtt{R}}^{\left(z\right)}
            +
        \left(yzx\right)
        \hat{U}_{\bm{r}}\,
        \Phi_{\mathtt{L}}^{\left(z\right)*}
        \Phi_{\mathtt{R}}^{\left(x\right)}
.
\end{align}